\subsection{Different probabilities: the rare type matters}
\label{sec:unequal-coupon-collector}

Now the $n$ types have probabilities $p_1,\ldots,p_n>0$, summing to
one; successive purchases are still independent. Let $T$ be the number
of purchases to collect all types, and write $p_A=\sum_{i\in A}p_i$.
If some required type has probability zero, completion is impossible
and $T=+\infty$; the finite formulas below assume positive probabilities.

The uniform stage argument needs revision. After collecting the set $S$,
the next-new-type probability is $1-\sum_{i\in S}p_i$. Its value depends
on the identities of the collected types, not just their number. For
probabilities $1/2,1/3,1/6$, a first coupon of type 1 leaves new-type
probability $1/2$, whereas a first coupon of type 3 leaves $5/6$.
There is no single geometric parameter determined only by the stage index.

\begin{proposition}[Exact unequal-probability collection formulas]
\label{prop:unequal-coupon-formulas}
For integers $m\ge0$,
\begin{align*}
\P(T>m)&=\sum_{\varnothing\ne A\subseteq\{1,\ldots,n\}}
(-1)^{|A|+1}(1-p_A)^m,\\
\E[T]&=\sum_{\varnothing\ne A\subseteq\{1,\ldots,n\}}
\frac{(-1)^{|A|+1}}{p_A}.
\end{align*}
In the tail formula at $m=0$, every $(1-p_A)^0$ means $1$, including
the full-set term: before any draw every specified set is missing.
\end{proposition}
\begin{proof}
For each type $i$, let $A_i$ be the event that it is missing after $m$
draws. A specified nonempty set $A$ is entirely missing with probability
$(1-p_A)^m$. The event $T>m$ is the union of the $A_i$;
inclusion--exclusion gives the tail formula. At zero draws the alternating
coefficients sum to one, as required.

Use $\E[T]=\sum_{m\ge0}\P(T>m)$. There are finitely many subsets,
and each geometric series converges because $p_A>0$. Summing each
series gives $\sum_{m\ge0}(1-p_A)^m=1/p_A$ and the stated mean.
This also proves that the mean is finite and completion is almost sure.
\end{proof}

\begin{example}[Three rarities]
For $(p_1,p_2,p_3)=(1/2,1/3,1/6)$,
\[
\E[T]=2+3+6-\left(\frac65+\frac32+2\right)+1
=\frac{73}{10}.
\]
For three equally likely types the mean is $3H_3=11/2$.
With only two types of probabilities $p,1-p$,
\[
\E[T]=\frac1p+\frac1{1-p}-1=\frac1{p(1-p)}-1\ge3,
\]
with equality exactly at $p=1/2$. As $p$ tends to zero the mean
diverges. A very common first type does not compensate for a very rare
last type.
\end{example}

The same calculation gives a second moment. Since
$T^2=\sum_{m=0}^{T-1}(2m+1)$,
\[
\E[T^2]=\sum_{\varnothing\ne A}(-1)^{|A|+1}
\frac{2-p_A}{p_A^2},\qquad
\Var(T)=\E[T^2]-(\E[T])^2.
\]
Indeed $\sum_{m\ge0}(2m+1)(1-p)^m=(2-p)/p^2$.
The sums involve finitely many convergent geometric series, so the
calculation is justified without assuming independent first-appearance times.

\subsection{Remaining waits and practical bounds}

Let $R(S)$ be the expected additional draws when precisely the types
in $S$ are still missing. Its boundary value is $R(\varnothing)=0$.
For nonempty $S$, first-step conditioning gives
\[
R(S)=1+(1-p_S)R(S)+\sum_{i\in S}p_iR(S\setminus\{i\}),
\]
and hence the positive recursion
\[
\boxed{R(S)=\frac{1+\sum_{i\in S}p_iR(S\setminus\{i\})}{p_S}.}
\]
For the three-rarity example, if types 2 and 3 are missing, the remaining
mean is $3+6-2=7$ draws. Knowing only that two types are missing is
insufficient. The subset formula for the mean applies to a missing set
as well, with the original, unnormalized probabilities: already collected
types still use up draws. The recursion provides another exact way to
compute these means, using only smaller missing sets.

Let $p_{\min}=\min_i p_i$. Waiting for the rarest type gives a lower
bound, while the maximum of the individual first-appearance times is
at most their sum. Thus
\[
\frac1{p_{\min}}\le\E[T]\le\sum_{i=1}^n\frac1{p_i}.
\]
For deadlines, linearity of missing-type indicators and the union bound give
\[
\E[M_m]=\sum_i(1-p_i)^m,\qquad
\P(T>m)\le\min\left\{1,\sum_i(1-p_i)^m\right\}
\le\min\{1,ne^{-mp_{\min}}\}.
\]
When $p_i=1/n$, the exact formulas reduce to the uniform formulas above;
group the subsets by cardinality. Section~\ref{sec:poissonization} gives
a second derivation and proves that the uniform choice minimizes the
expected completion time among all positive probability vectors.
